\section{引言}
近地轨道是对地观测、空间实验和航天运输的重要运行区域。入轨不仅要求给定高度，还要求位置与速度同时满足轨道条件。齐奥尔科夫斯基方程描述理想推进能力与质量比的关系\cite{tsiolkovsky1903,sutton2016}，但不能直接给出包含引力、阻力和有限推力转向的实际轨迹。

控制参数化与动力学积分可将入轨问题转化为有限维非线性优化。PSO通过个体与群体经验搜索\cite{kennedy1995pso}，GA通过选择与重组传播候选信息\cite{holland1992}，SA利用温度相关概率接受规则探索非凸目标\cite{kirkpatrick1983}。表现依赖边界、初始点及计算预算，单次最好结果不足以支持方法间比较；附加局部精修还会混合不同算法的贡献。

离线可行轨迹也不保证在新扰动下准确开环执行。MPC通过反馈、有限时域预测和约束优化修正输入\cite{mayne2000mpc,borrelli2017}，适合协调参考附近的状态误差与控制边界。

本文在统一动力学和两段燃烧参数化下比较无局部精修的PSO、GA和SA，区分搜索评价与高精度复核。随后以纯PSO轨迹构造时变线性化MPC，检验同一实际扰动下的开闭环差异。结论限于单级质点算例及所用实现，不作算法一般性排序或工程适用性声明。
